\documentclass[11pt,a4paper]{article}
\usepackage[margin=25mm]{geometry}
\usepackage{amsmath,amssymb,siunitx,hyperref,listings,microtype}
\hypersetup{colorlinks=true,linkcolor=blue,urlcolor=blue,pdftitle={Report 6 - Shaft Model Family for OpenHPL/OpenHPLjl},pdfauthor={Madhusudhan Pandey}}
\lstset{basicstyle=\ttfamily\small,breaklines=true,frame=single}
\title{Report 6 - Shaft Model Family for OpenHPL/OpenHPLjl}
\author{Madhusudhan Pandey}
\date{19 September 2026}
\begin{document}
\maketitle
\begin{abstract}
This report introduces the shaft model family as the mechanical bridge between turbine torque and generator torque. The first implementation is a lumped single-inertia shaft with viscous damping and an acausal rotational connector.
\end{abstract}
\section{Updated Context}
The shaft transfers turbine mechanical power to the generator:
\[
\boxed{\text{Turbine}\rightarrow\text{Shaft}\rightarrow\text{Generator}}.
\]
Its first role in OpenHPLjl is to represent rotational inertia and speed dynamics.

\section{Generalized Concept Model}
The shaft receives drive and load torques and exposes one common angular speed:
\[
\tau_{drive},\tau_{load}\rightarrow J\rightarrow\omega.
\]

\section{Other Models Available in the Literature}
Useful model levels are
\[
\boxed{\text{rigid massless shaft}\rightarrow\text{single inertia}\rightarrow\text{two-mass torsional shaft}\rightarrow\text{distributed flexible shaft}}.
\]

\section{Mathematics Involved}
For the lumped shaft,
\[
\boxed{J\dot\omega=\tau_{drive}+\tau_{load}-D(\omega-\omega_0)}.
\]
Mechanical power is
\[
P_m=\tau\omega.
\]

\section{OpenHPL-Specific Modeling Interpretation}
OpenHPL includes rotational mechanical connections between turbine and generator. OpenHPLjl mirrors that physical interface with an acausal rotational port containing angular speed and flow-conserved torque.

\section{Implementation in Julia ModelingToolkit / SciML}
\begin{lstlisting}
@connector RotationalPort begin
    omega(t)
    tau(t), [connect = Flow]
end

@component function LumpedShaft(; name,
    J=1e5, damping=0.0, omega0=2*pi*50)
    ...
    J * D(omega) ~ drive.tau + load.tau -
                   damping*(omega-omega0)
end
\end{lstlisting}

\section{Model Assembly and Connections}
The shaft is intended to connect the turbine mechanical port to the generator mechanical port. The next package step will give the turbine an explicit rotational connector.

\section{Numerical Experiment / Validation}
For
\[
J=100,\quad\tau_{drive}=20,\quad\tau_{load}=-10,
\]
the acceleration is
\[
\dot\omega=\frac{10}{100}=0.1~\si{\radian\per\second\squared}.
\]
Starting at $100~\si{\radian\per\second}$,
\[
\boxed{\omega(2)=100.2~\si{\radian\per\second}}.
\]

\section{Expected Outputs}
Key outputs are
\[
\omega(t),\quad\dot\omega(t),\quad\tau_{drive}(t),\quad\tau_{load}(t).
\]

\section{Engineering Interpretation}
The shaft introduces mechanical inertia, the central storage mechanism governing electromechanical frequency dynamics.

\section*{Conclusion}
The first shaft model is intentionally minimal. Next steps are a two-mass torsional model and explicit coupling to turbine torque and generator electromagnetic torque.
\end{document}
